% Generated by scripts/make_paper_tables.py -- do not edit.
Every result in the paper is measured in two worlds that differ on every axis that
could plausibly matter, so that a finding which depends on the construction can be told from one
that does not. \textsc{kitchen} is the important one: its evidence is \emph{non-additive}, so
losing a cue can be catastrophic rather than gradual, which is the structure most likely to break
a mechanism map built on gradual evidence loss. Cue visibilities are deliberately spread rather
than uniform --- if every cue were equally visible the observability mask would carry no
information and every scene would sit at the same irreducible floor, which is the degenerate case
a routing experiment must avoid.

\begin{table}[h]
\centering
\small
\caption{The two worlds. Only the second admits interactions between cues, and the instrument's
checks are world-independent by construction rather than by configuration.}
\begin{tabular}{lcc}
\toprule
& \textsc{junction} & \textsc{kitchen} \\
\midrule
outcomes & 4 & 3 \\
cues & 4 & 5 \\
cue alphabets & uniformly ternary & mixed binary and ternary \\
symmetry & left/right mirrored & none \\
evidence & \textbf{additive} (linear in the cues) & \textbf{non-additive} (pairwise) \\
joint cue states & \JointStates & \KitchenStates \\
irreducible floor at the default $\tau$ & \FloorJ{} nats & \FloorK{} nats \\
validation checks passed & \NumChecks/\NumChecks & \NumChecks/\NumChecks \\
\bottomrule
\end{tabular}
\end{table}
